\documentclass[addpoints,12pt]{exam}

\usepackage{geometry}
 \geometry{
 a4paper,
 total={170mm,260mm},
 left=20mm,
 top=15mm,
 }

\usepackage{amsmath,amssymb,amsthm}
\DeclareFontEncoding{T1}
\fontencoding{T1}
\usepackage[T2A]{fontenc} 
\usepackage[utf8]{inputenc}
\usepackage[russian]{babel}
\usepackage{graphics,graphicx}

\usepackage{verbatim}
\usepackage{fancyvrb}  

\usepackage{listings}
\usepackage{color}

\renewcommand{\solutiontitle}{\noindent\textbf{Решение:}\par\noindent}

\definecolor{codegray}{gray}{0.95}
\definecolor{keywordcolor}{RGB}{0,0,255}
\definecolor{commentcolor}{RGB}{0,128,0}

\lstset{
    language=Matlab,
    basicstyle=\ttfamily\footnotesize,
    backgroundcolor=\color{codegray},
    keywordstyle=\color{keywordcolor},
    commentstyle=\color{commentcolor},
    numbers=left,
    numberstyle=\tiny\color{gray},
    numbersep=5pt,
    frame=single,
    tabsize=2,
    breaklines=true,
    captionpos=b,
    escapeinside={\%*}{*)}
}

\renewcommand{\solutiontitle}{\noindent\textbf{Решение:}\par\noindent}

\definecolor{gray}{gray}{0.5}

\lstset{basicstyle=\ttfamily,
    language=Matlab,    
    keepspaces=true,
    extendedchars=\true,
    identifierstyle={ \bf \color{black} },
    showstringspaces=false,
    numbers=left,%
    numbersep=7pt,  
    emph=[1]{case,switch,otherwise,nonlocal,as,yield,zip,print},    
    frame=single,    
    xleftmargin=0.3cm,
    frame=l,framesep=4pt,framerule=0.5pt,
    numberstyle={\scriptsize \color{gray}}
}


\pagestyle{headandfoot}
\runningheadrule
\firstpageheader{ИнтМатПак}{MATLAB 1}{17 февраля 2026}
\runningheader{ИнтМатПак}
{MATLAB 1, Стр. \thepage\ из \numpages}
{17 февраля 2026}
\firstpagefooter{}{}{}
\runningfooter{}{}{}
\pointname{} 
\boxedpoints
%\pointsinmargin
\pointsinrightmargin


\begin{document}

\makebox[\textwidth]{Ф. И. О. \enspace \hrulefill  группа: \enspace\hrulefill}
\vspace{0.3cm}

\begin{questions}
 
\question[1] Чему будет равна переменная B в результате выполнения кода?
\begin{lstlisting}
M = reshape(1:16, 4, 4)';
B = M(2:3, [1,4]);
\end{lstlisting}
\begin{checkboxes}
\choice \texttt{[2 3; 6 7]} 
\choice \texttt{[5 8; 6 7]}
\CorrectChoice \texttt{[5 8; 9 12]}
\choice \texttt{[6 7; 10 11]}
\end{checkboxes}

\question[1] Что будет содержаться в переменной \texttt{result}?
\begin{lstlisting}
A = [1, 0, 3; 4, 5, 0; 0, 8, 9];
result = A(A > 2);
\end{lstlisting}
\begin{checkboxes}
\choice \texttt{[1; 0; 3]}
\CorrectChoice \texttt{[4; 5; 8; 3; 9]}
\choice \texttt{[4; 5; 8; 9]}
\choice \texttt{[3; 4; 5; 8; 9]}
\end{checkboxes}

\question[1] Чему будет равна переменная b?
\begin{lstlisting}
a = [7,3,6,8,2,4];
b = reshape(a,2,3);
\end{lstlisting}
\begin{checkboxes}
\choice  [7, 3, 8; 6, 2, 4]
\CorrectChoice [7, 6, 2; 3, 8, 4]
\choice  [7, 3; 8, 6; 2, 4]
\choice  [7, 6; 2, 3; 8, 4]
\end{checkboxes}

\question[1] Есть матрица \texttt{temperatures} размером 30×24, представляющая температуру (в градусах Цельсия) каждый час в течение 30 дней (строка — день, столбец — час). Напишите код, который определяет среднюю температуру для каждого дня (результат — вектор-столбец 30×1).
\begin{solutionorlines}[2cm]
\end{solutionorlines}
 
\question[1] Чему будет равна переменная i в результате выполнения кода?
\begin{lstlisting}
a = [7,3,6,4];
[~,i] = sort(a);
\end{lstlisting}
\begin{checkboxes}
\CorrectChoice \texttt{[2,4,3,1]}
\choice \texttt{[1,3,4,2]}
\choice \texttt{[3,4,6,7]}
\choice \texttt{[7,6,4,3]}
\end{checkboxes}

\question[1] Дан вектор \texttt{v = 1:12}. Напишите код, который преобразует этот вектор в матрицу 3×4, заполняя её построчно.
\begin{solutionorlines}[2cm]
\end{solutionorlines}

\question[1] Чему будет равна переменная a?
\begin{lstlisting}
a = sum([7,6,3,9].*[0,1,0,1])
\end{lstlisting}    
\begin{checkboxes}
\choice 7;
\choice 9;
\CorrectChoice 15;
\choice выведется сообщение об ошибке.
\end{checkboxes}

\question[1] Координаты двух точек в пространстве заданы матрицами-строками \texttt{A = [1,2,0]; B = [3,1,1]}. Напишите код, который определяет расстояние между этими двумя точками.
\begin{solutionorlines}[1cm]
\end{solutionorlines}

\question[1] 
Дана матрица \texttt{X} размером 50x50, заполненная нормально-распределенными случайными числами (мат. ожидание равно 20, стандартное отклонение равно 10): 
\begin{lstlisting}
X = randn(50, 50) * 10 + 20; 
\end{lstlisting}
Напишите код, который извлекает <<периметр>> матрицы X (первая и последняя строки, первый и последний столбцы) в линейный вектор \texttt{perimeter\_vec}, идущий по часовой стрелке, начиная с элемента (1,1).
\begin{solutionorlines}[3cm]
    \texttt{perimeter\_vec} = 
\end{solutionorlines}

\question[1] 
Напишите код, который в последнем столбце матрицы \texttt{X} из предыдущего задания все значения, превышающие 30, увеличивает в два раза (должна изменится матрица X).
\begin{solutionorlines}[2cm]
\end{solutionorlines}

\end{questions}

\hqword{Вопрос:}
\hpword{Баллы:}
\hsword{Ответ:}
\htword{Сумма}

\vfill

\begin{center}
\gradetable[h][questions]
\end{center}

\end{document}